\documentclass{article}

% Language setting
% Replace `english' with e.g. `spanish' to change the document language
\usepackage[english]{babel}

% Set page size and margins
% Replace `letterpaper' with `a4paper' for UK/EU standard size
\usepackage[letterpaper,top=2cm,bottom=2cm,left=3cm,right=3cm,marginparwidth=1.75cm]{geometry}

% Useful packages
\usepackage{amsmath}
\usepackage{graphicx}
\usepackage[colorlinks=true, allcolors=blue]{hyperref}
\usepackage{pgfgantt}
\usepackage{enumerate}
\usepackage{enumitem}


\begin{document}

\section{Cronograma}

Apresentação da seção.

\begin{figure}[h]
	\caption{Diagrama de Gantt ilustrando o cronograma das atividades planejadas ao longo do projeto, com a linha do tempo distribuída em intervalos quinzenais}
	
	\begin{ganttchart}[
		y unit title=0.7cm,
		y unit chart=0.5cm,
		x unit=0.8cm,
		vgrid,hgrid,
		title height=1,
		title/.style={fill=none},
		title label font=\bfseries\footnotesize,
		bar/.style={fill=cyan},
		bar height=0.7,
		progress label text={},
		group right shift=0,
		group top shift=0.7,
		group height=.3,
		group peaks width={0.2},
		inline]{1}{20}
		
		\gantttitle{2024}{10}
		\gantttitle{2025}{10}\\
		
		\gantttitle{Ago}{2}
		\gantttitle{Set}{2}
		\gantttitle{Out}{2}
		\gantttitle{Nov}{2}
		\gantttitle{Dez}{2}
		\gantttitle{Fev}{2}
		\gantttitle{Mar}{2}
		\gantttitle{Abr}{2}
		\gantttitle{Mai}{2}
		\gantttitle{Jun}{2}\\
		
		\ganttbar[progress=100,inline=false]{1)}{1}{1}   \\ % 1
		\ganttbar[progress=100,inline=false]{2)}{2}{2}   \\ % 2
		\ganttbar[progress=100,inline=false]{3)}{2}{3}   \\ % 3
		\ganttbar[progress=100,inline=false]{4)}{3}{4}   \\ % 4
		\ganttbar[progress=60,inline=false]{5)}{3}{6}    \\ % 5
		\ganttbar[progress=0,inline=false]{6)}{5}{8}     \\ % 6
		\ganttbar[progress=0,inline=false]{7)}{7}{8}     \\ % 7
		\ganttbar[progress=0,inline=false]{8)}{8}{11}    \\ % 8
        \ganttbar[progress=0,inline=false]{9)}{9}{9}     \\ % 9
		\ganttbar[progress=0,inline=false]{10)}{11}{14}  \\ % 10
		\ganttbar[progress=0,inline=false]{11)}{12}{15}  \\ % 11
		\ganttbar[progress=0,inline=false]{12)}{12}{16}  \\ % 12
		\ganttbar[progress=0,inline=false]{13)}{15}{18}  \\ % 13
		\ganttbar[progress=0,inline=false]{14)}{16}{18}  \\ %14
		\ganttbar[progress=0,inline=false]{15)}{15}{19}  \\ % 15
		\ganttbar[progress=0,inline=false]{16)}{20}{20}     % 16
	\end{ganttchart}
	
	\begin{enumerate}[nosep, label=\arabic*)]
		\item Definição de Orientador. % 1
  
		\item Escolha do Tema. % 2
		
		\item Definição da Pergunta de Pesquisa. % 3
		
		\item Delimitação do Tema e Objetivos. % 4
		
		\item Escrita de Materiais e Métodos. % 5
		
        \item Implementação Ambiente de Teste % 6
        
        \item Validação do Ambiente de Experimento % 7
        
        \item Preparação para ações de simulação ou coleta de dados % 8

        \item Apresentação da proposta de trabalho % 9
        
        \item Coleta de Dados % 10
        
        \item Análise dos dados % 11
        
        \item Redação do texto dos experimentos % 12
        
        \item Redação do texto de Análise, Discussão dos resultados e Considerações Finais % 13
        
        \item Finalização do texto % 14
        
        \item Preparação para apresentação % 15
        
        \item Revisão Final % 16

	\end{enumerate}
	
	\label{fig:gantt-schedule}
	
\end{figure}


\pagebreak

\end{document}